% TOP_PROBLEM: {"problemId": "problem.the-abc-conjecture-of-masser-and-oesterle", "canonicalTitle": "The abc Conjecture of Masser and Oesterle", "releaseRank": 11, "primaryDomain": "number_theory_arithmetic_geometry", "primaryDomainLabel": "Number theory & arithmetic geometry", "displayStatus": "open_disputed_claim", "releaseStatus": "open_disputed_claim"}
% !TeX program = lualatex
% Definition and sources only; this file is not an LLM solution attempt.
\documentclass[11pt]{article}
\usepackage[margin=25mm]{geometry}
\usepackage{fontspec}
\setmainfont{Latin Modern Roman}
\usepackage{amsmath,amssymb,mathtools}
\usepackage{unicode-math}
\setmathfont{Latin Modern Math}
\usepackage{xurl,hyperref}
\hypersetup{colorlinks=true,urlcolor=blue}
\setcounter{secnumdepth}{0}
\setlength{\parindent}{0pt}
\setlength{\parskip}{5pt}
\setlength{\emergencystretch}{3em}
\begin{document}
\section*{\texorpdfstring{The abc Conjecture of Masser and Oesterle}{The abc Conjecture of Masser and Oesterle}}\label{11}
\subsection{Definitions and mathematical statement}
For a positive integer $n$, its radical is $\operatorname{rad}(n)=\prod_{p\mid n}p$, with $\operatorname{rad}(1)=1$. A triple $a,b,c>0$ with $a+b=c$ and $\gcd(a,b)=1$ is automatically pairwise coprime. The conjecture is
\[
 \forall\varepsilon>0\ \exists C_\varepsilon>0\ \forall a,b,c\in{\mathbb{Z}}_{>0},\qquad
 \begin{cases}a+b=c,\\\gcd(a,b)=1\end{cases}
 \Longrightarrow c<C_\varepsilon\operatorname{rad}(abc)^{1+\varepsilon}.
\]
The constant depends on $\varepsilon$ only. A counterexample must defeat every constant for some fixed positive $\varepsilon$; one unusually large triple does not suffice.
\par\smallskip\textit{Formulation and concept references:} \href{https://link.springer.com/article/10.1007/s00025-026-02616-5}{[S1]}.

\subsection{Short English statement}
Can the size of a coprime additive triple always be controlled, up to any small exponent loss, by its distinct prime factors?
\subsection{Sources}
\begin{itemize}
\item[S1]Distance Between Cubics and Rationals, Results in Mathematics (2026), statement of the abc conjecture.
\url{https://link.springer.com/article/10.1007/s00025-026-02616-5}.
\textit{Formulation source.}
\item[S2]Construction of Arithmetic Teichmuller Spaces IV: Proof of the abc-conjecture — Kirti Joshi.
\url{https://arxiv.org/abs/2403.10430}.
\textit{Further reference; not a proof endorsement.}
\item[S3]Why abc is still a conjecture — Peter Scholze and Jakob Stix, July 16, 2018.
\url{https://www.math.uni-bonn.de/people/scholze/WhyABCisStillaConjecture.pdf}.
\textit{Further reference; not a proof endorsement.}
\item[S4]The abc Conjecture Revisited — Patrick Letendre, July 2026.
\url{https://arxiv.org/abs/2607.07641}.
\textit{Further reference; not a proof endorsement.}
\end{itemize}
\end{document}
